\item On the signed arc digraph of a line, a rule that adds over end vertices a function of the modulus of the summed gains is consistent under continuation only if that function is the square (Theorem~\ref{thm:quad}); the same holds on every finite connected signed digraph that is norm-preserving of degree two (Theorem~\ref{thm:deg2}). On any norm-preserving gain digraph with a merging pair of walks, the only consistent power of the modulus is the second (Theorem~\ref{thm:power}). On a norm-preserving carrier of degree two with gains in the $q$-th roots of unity, $q$ even, on which every gain reaches both vertices of a block in abundance, the consistent rules of this form are exactly the combinations of the squared moduli in the conjugate characters; they form a space of dimension $\varphi(q)/2$ for $q>2$ and of dimension one for $q=2$, and the rule is unique up to a factor exactly for $q\in\{2,4,6\}$ (Theorem~\ref{thm:conj}). Such a rule is non-negative exactly when its weights are, and if it is multiplicative on products of carriers it is the squared modulus in one conjugate character (Proposition~\ref{prop:compo}). If the gain of a walk is itself recorded, so that the carrier is the signed lift with one sheet for each pair of opposite gains, the rule is unique for every $q$: it is the square of the integer amplitude, the mean of the squared moduli over the characters in which $-1$ acts as $-1$ (Corollary~\ref{cor:lift}). Without any choice of records, for any finite abelian group of gains and any degree, and under a hypothesis on the quotients of gains that reach two vertices with a common successor, the consistent switching-invariant pair rules of \cite{g8} are exactly the real combinations of the squared-norm rules of the characters for which the step is norm-preserving (Theorem~\ref{thm:kernel}). In degree two these are the characters in which $-1$ acts as $-1$, the kernel changes sign under $-1$, and for the $q$-th roots of unity the space has dimension $\lceil q/4\rceil$ and contains both families (Corollary~\ref{cor:kernel2}). For blocks of degree three and four with gains of order $3$, $6$ or $4$, with disjoint successor sets and abundant gains, a consistent rule that adds a function of the modulus over end vertices is again the square (Proposition~\ref{prop:deg34}); for blocks $\rho_z\kappa_yw^{zy}$ of every degree $d$ dividing $q$, under the same conditions, the consistent rules of that form are exactly the combinations of the conjugate squared moduli, as in degree two (Theorem~\ref{thm:fourier}). Finally the form of the blocks is not needed at all: on every norm-preserving carrier, of any degree, in which two vertices have a common successor and the events at them are rich enough, the consistent rules of that form are exactly the combinations of the conjugate squared moduli (Theorem~\ref{thm:vsl}, Corollary~\ref{cor:vsl}). On a finite carrier whose derived digraph is strongly connected no hypothesis on events is needed: the order that enters is that of the gains that can meet at one length (Proposition~\ref{prop:abund}, Corollary~\ref{cor:abund}). On carriers with a full block, that is, $d$ vertices with one common set of $d$ successors and a matrix of gains $B$ with $B^*B=d$ for the faithful character, and with abundant gains, the modulus is not needed either; every finite norm-preserving digraph of degree two consists of such blocks. A consistent rule that adds over end vertices any function of the summed gains that vanishes at zero is a quadratic form, and if it is unchanged by switching it is a combination of the conjugate squared moduli (Theorem~\ref{thm:block}, Corollary~\ref{cor:block}). Nor is the sum: for any $d\times d$ matrix of gains in a finite abelian group, $d\ge2$, a function of the multiset of gains that satisfies the block law is an additive function plus a quadratic form, and if the entries generate the group it is a combination of the count and of squared moduli of gains summed through characters; on a carrier with such a block and abundant gains this holds for the consistent rules (Theorem~\ref{thm:multiset}, Corollaries~\ref{cor:multiset}, \ref{cor:multicar}). Nor need the function be the same at every end vertex, or the carrier have a block: on a finite strongly connected gain digraph in which every vertex has $d\ge2$ outgoing arcs, without parallel arcs and with abundant gains, a consistent rule that adds over end vertices a function of the multiset of gains that may depend on the vertex, and vanishes on the empty multiset, is a combination of the count and of the squares $|\sigma_\psi|^2$ for the characters $\psi$ in which the step is norm-preserving, with the same coefficients at every vertex (Theorem~\ref{thm:vertexlaw}, Corollary~\ref{cor:vertexlaw}). The abundance of every gain at every vertex at one length fails on periodic carriers such as a ring with an even number of sites; it is not needed: if the derived digraph on pairs (vertex, gain) is strongly connected, of any period, the consistent rules of this kind are, as functions on events, exactly $a|S|+\sum c\,\lVert v^\psi_S\rVert^2$ over those characters, with constants that depend neither on the vertex nor on the length, and characters whose restrictions to the subgroup of quotients of gains that can meet at one vertex at one length are equal or inverse give the same rule (Theorem~\ref{thm:periodic}, Corollary~\ref{cor:periodic}). If such a consistent rule is in addition multiplicative on the product events of the carrier with itself, at all large lengths, it is zero, the count, or the squared norm for one of these classes of characters: mixtures of the count and a square, and of two squares, are excluded (Corollary~\ref{cor:composition}). For a non-negative consistent rule on such a carrier, normalized so that one walk has weight one, with count coefficient $a$, the weight of the walks that return to the starting vertex, divided by the number of all walks, is at most $1-a(1-1/d)$ at every positive length and comes arbitrarily close to $1-a$ or more; so a requirement that it come arbitrarily close to one at positive lengths excludes the count (Corollary~\ref{cor:return}; the idea is Director~G's \cite{dgret}). The return weight is computed with the rule's own weights, so this requirement is a condition on the rule, equivalent to the vanishing of its count part, and not an independent fact about the system.
Director~G observed \cite{dgg1} that Theorem~\ref{thm:quad} needs far fewer events when $\phi$ is taken as a function of the signed sum: if only the events of all walks that end at one site are used (both rails together), the same $a+1$ walks give the same recursion, the factor $c=2$ follows from the site at the front and need not be assumed, and the sites at the other edge give $\phi(-n)=\phi(n)$. A computation made for the review of that note shows what is traded: with these events alone, a function of the signed sum is forced to be the square, a function of the number of walks is forced to be the count, and a function of both is not forced at all (the space of solutions has dimension $37$ at length $12$); if the same demand is made for every launch with integer entries and not only for a single walk, a function $F$ of the numbers $n_+$, $n_-$ of walks of either gain with $F(0,0)=0$ is $a(n_++n_-)+c(n_+-n_-)^2$ \cite{comp4s}. These statements are records of a review and are not among the statements of this paper.

So if the weight of returning to the starting vertex comes arbitrarily close to one at positive times, which Director~G argues it must when the complete system has a compact state space of finite dimension and a reversible evolution \cite{dgret} (that argument is not part of this paper), the rule has no count part: it is a non-negative combination of squares, and if it is also multiplicative in the sense of Corollary~\ref{cor:composition} it is the squared norm for one class of characters. The count is the rule of a random walk in which each of the $d$ arcs is taken with equal chance; under it the return weight never exceeds $1/d$ after the start. The return weight is defined by the weights of the rule itself. By (3) the requirement that it come close to one is therefore equivalent to $a=0$: it is a condition on the rule and cannot serve as an independent reason for the absence of the count. An earlier edition called it a premise about the system; that was wrong, and a stipulated candidate model built on it was withdrawn for this reason \cite{comp4s}. An independent premise has to be about measured frequencies. Director~G gave two on the signed ring: one port reached by exactly two walks of opposite gain has weight $a/4$ at length three, so that it is dark exactly when $a=0$; and the port $(1,R)$ has weight $\tfrac12$ at length one and $\tfrac12-\tfrac a4$ at length three, so that it is no dimmer at length three exactly when $a=0$ \cite{dgz1}. The signs assumed in the corollary need not be assumed either: if the rule is non-negative on all events, of unbounded length, then $a\ge0$ and every $c_{\psi'}\ge0$, the classes being distinct as restrictions to $C_0$; on events with at most $M$ walks of each gain at a vertex the rule $(1-c)(r+s)+c(r-s)^2$ with $c=-1/(M-1)$ is still non-negative, so bounded tests do not give the signs (Director~G's proof \cite{dgz1}; a challenger found no failing step and, among $506$ rules with a negative coefficient, a negative event for each, $107$ of them only beyond multiplicity three). It uses the starting vertex, positivity and the normalization; the strong connectivity of the derived digraph enters only through the form of the rule. The bound (1) is attained (two vertices with the step matrix $\bigl(\begin{smallmatrix}1&1\\1&-1\end{smallmatrix}\bigr)$, $t=2$), and it fails when parallel arcs are allowed. On $11$ carriers (two vertices over the group of order two, $B_q$ for $q=4,6,8,12$, rings of $3$, $4$, $5$, $6$ and $8$ sites, a Fourier block of degree three) the bound (1) holds at every length up to $4000$ for $a=\tfrac14,\tfrac12,1$ and the coherent part returns above $0.99$ within that length; this depends on the carriers chosen and on the horizon: on signed rings of $9$, $11$ and $12$ sites the first return above $0.99$ is at length $95946$, beyond $100000$ and at $6416$, and the return weight of the square equals one at some positive length only on the rings of $2$, $4$ and $8$ sites among those of $2$ to $40$ sites; twelve random carriers did not satisfy the hypothesis that some character is norm-preserving and are listed \cite{comp4s}. A first run of that program stopped while writing its result and is kept.
Corollary~\ref{cor:composition} and its proof were written by the author; a challenger found no failing step and, on four carriers, by enumeration of all pairs of events of length three (up to $65{,}536$ pairs), that exactly the rules named are multiplicative, and that at length two some mixtures are multiplicative as well, which is why the hypothesis asks for large lengths; the algebraic step $e_ie_j=\delta_{ij}e_i$ is checked in Lean~4 with Mathlib over any commutative ring, and its consequence that each $e_i$ is $0$ or $1$ over a domain \cite{comp4s}. Carriers with $C_0\ne C$ were not computed for this corollary. The corollary concerns the product of a carrier with itself in which a pair of walks is recorded by the product of its gains at a pair of vertices. It says nothing about composites whose parts keep records of their own: Director~G gave a rule on the signed ring with a private binary record for each part that is multiplicative on all product events and is not a function of the gains at the end vertices \cite{dgz1}.
